% Auto-converted from Optimizing_Gestural_Efficiency.docx via pandoc
% Source section: Algorithm and Design
Concepts
\chapter{Algorithm and Design Concepts}\label{algorithm-and-design-concepts}

\section{The Algorithm}\label{the-algorithm}
INPUTS

Vectors

Name Elements Ordered By Input/Calculated

vLexeme: Vector of Lexemes ordered by Corpus Frequency

vLexon: Vector of Lexon ordered by Corpus Frequency

vQwerty: Vector of Qwerty symbols ordered by `User Convenience'

\section{Prefix Match versus Exact
Match}\label{prefix-match-versus-exact-match}

\section{Prefixes are GridNames}\label{prefixes-are-gridnames}
\section{User-Convenience
Ordering}\label{user-convenience-ordering}

\section{Proof of Correctness}\label{proof-of-correctness}
\section{Proof of Termination}\label{proof-of-termination}
Alphabetical(TouchTyping), HomeRow(Effort), KbdPosition(Location)

\section{Reasonability of the Measure for Gestural
Efficiency}\label{reasonability-of-the-measure-for-gestural-efficiency}

\section{Minimizing Interference with Target
Apps}\label{minimizing-interference-with-target-apps}

\emph{\textbf{µLex}}'s \emph{raison d'être} is enabling non-QWERTY
lexical input GFAs (see section 2.7) for \textbf{\uline{other}} apps,
such as word processors and email managers. In OS terms, other apps for
which \emph{\textbf{µLex}} provides supplemental GFAs keep
\emph{\textbf{focus}}. Since \emph{\textbf{µLex}} apps work in
conjunction with focus apps, \emph{\textbf{µLex}} GFAs must not
interfere with focus app GFAs. To minimize interference:

\begin{itemize}
\item
  \textbf{DC01}: \emph{\textbf{µLex}} apps are not running by default.
\item
  \textbf{DC02}: Potentially interfering gestures are only active when
  \emph{\textbf{µLex}} is toggled on.
\item
  \textbf{DC03}: Toggling \emph{\textbf{µLex}} is only possible with
  gestures unused by focus apps.
\end{itemize}

\section{\texorpdfstring{ Non-Interference via Activation and
Toggle (DC01,
DC04)}{ Non-Interference via Activation and Toggle (DC01, DC04)}}\label{non-interference-via-activation-and-toggle-dc01-dc04}

\section{\texorpdfstring{\emph{\textbf{µLex}} Supplemental Displays
Prompt
Selections}{µLex Supplemental Displays Prompt Selections}}\label{uxb5lex-supplemental-displays-prompt-selections}

When \textbf{\emph{µ}短语} is active, cells indicate which lexemes can
be input with gestures. Blue letters indicate keyboard gestures and cell
positions indicate mouse gestures. If chordsets are used, the
\textbf{\emph{µ}短语} display is supplemented with chordmaps. Chordmaps
indicates GFAs mapping finger combinations to lexemes. So, when
\textbf{\emph{µ}短语} is toggled on, the lexeme \textbf{为} (wèi
`because of') can be input three ways:

\begin{itemize}
\item
  via keyboard, by actuating lowercase-z.
\item
  via mouse, by clicking the containing region.
\item
  via chordset, by chording thumb, ring, and little fingers.
\end{itemize}

\begin{figure}[htbp]
\centering
\includegraphics[width=6.5in,height=2.48403in]{media/media/image14.emf}
\caption{Chordmaps Indicate Chordset GFAs}
\label{fig:auto:1}
\end{figure}


\section{\texorpdfstring{\emph{\textbf{µLex}} Supplemental Actuation
Methods Enable Alternative
Selections}{µLex Supplemental Actuation Methods Enable Alternative Selections}}\label{uxb5lex-supplemental-actuation-methods-enable-alternative-selections}

\emph{\textbf{µLex}} provides special actuation methods for mice,
keyboards, and chordsets to enable users seamlessly to alternate between
focus app GFAs and \emph{\textbf{µLex}} GFAs.

\emph{\textbf{µLex}} apps coexist with and supplement apps such as word
processors and spreadsheets to allow users easily to input non-QWERTY
lexemes.

The white and cyan \emph{\textbf{cells}} in the initial
\textbf{\emph{µ}短语} display \emph{\textbf{grid}} enables users to
input the 30 most common lexemes (accounting for 33\% of lexical input)
with a single \textbf{\uline{primary}} gesture.

For mice, users move the cursor over a selection

\section{DC01: Supplement Existing Apps and
Devices}\label{dc01-supplement-existing-apps-and-devices}

To not interfere with keyboard accelerators and other keyboard GFAs,
\emph{\textbf{µLex}} software supplements devices with an additional
alternative actuation: \emph{\textbf{long press}} (sustaining actuation
for more than the length of time for a normal key press). One use of
long press is to toggle \emph{\textbf{µLex}} software availability. For
example, to alternate between English and Chinese lexical input, users
long press the \textbf{/} key, which shows or hides the supplemental
\emph{\textbf{grid}} of \emph{\textbf{cells}} shown in Figure 14.

When the display is toggled off (hidden), the keyboard behaves normally;
letter, number, and punctuation keys input their characters to the
currently active app. When the display is toggled on (visible), keys
(parenthesized blue) are used for selecting the indicated Chinese
lexemes into the currently active app. For example, Figure 14 indicates
actuating the 1 key inputs \textbf{的} (the most frequently used Chinese
lexeme). Input of Chinese lexemes not shown in Figure 14 will be
discussed later.

(Figure 18). In the display (Figure 18, right) the gestures are 31
different combinations of actuated fingers. Each combination is
associated with a \emph{\textbf{chordmap}} ((Figure 18, center)
indicating the fingers to be pressed and (simultaneously) released. The
display grid associates chordmaps with Chinese characters pronounced
\emph{\textbf{shì}}. To input characters, users actuate chords (Figure
18, left/green) indicated by chordmaps.

\begin{figure}[htbp]
\centering
\includegraphics[width=6.5in,height=2.26736in]{media/media/image15.emf}
\caption{Sample \emph{\textbf{µLex}} Display for Chinese Input}
\label{fig:auto:2}
\end{figure}


What distinguishes the keyboard from other devices is the
\textbf{\uline{marks}} on the keys. The marks constitute a display of
user choices\emph{\textbf{.}} For devices that don't have marks (e.g.,
mice, chordsets) the apps themselves present GUI displays for their
usage. For mice, the display is often the GUI itself, For
non-positional, unmarked devices an additional display must be provided
to For devices other than the keyboard it is trivial to supplement
existing apps with a display.

Every GUI-based app that takes lexical input (e.g., email, word
processing) is a potential \emph{\textbf{client}} for
\emph{\textbf{µLex}} software. Many clients have gestures

so, \emph{\textbf{µLex}} must operate cooperatively with its clients. In
particular \emph{\textbf{µLex}} gestures must be different from gestures
already used by the clients.

\section{Long Press -- A Supplemental Alternative
Actuation}\label{long-press-a-supplemental-alternative-actuation}

As a springboard example, \textbf{\emph{µ}Letras} is a
\emph{\textbf{µLex}} app that enables users to input non-QWERTY Spanish
characters from a keyboard with a gestural efficiency of 1. The
essential problem addressed by \textbf{\emph{µ}Letras} is that many
characters used in Spanish documents are not available on the QWERTY
keyboard. Solving the problem is simplified by the observation that
effectively all non-QWERTY characters are in 1-to-1 correspondence with
QWERTY marks.

The solution is made harder by the variety and multiplicity of keyboard
accelerators in software apps via alternative actuations of many keys.
For example
\textbf{\textless ctrl\textgreater-\textless a\textgreater{}} in
Microsoft word invokes the command `select all text in document'; so,
\textbf{\textless ctrl\textgreater-\textless a\textgreater{}} cannot be
used to insert characters. Similarly, alternative actuations using
combinations of \textless alt\textgreater, \textless esc\textgreater,
\textless ctrl\textgreater, and \textless shift\textgreater{} with other
QWERTY keys makes them infeasible for inserting non-QWERTY Unicode
characters.

The \textbf{\emph{µ}Letras} lexicon comprises phrases composed of the
QWERTY marks supplemented by the characters in Figure 19. Inspection of
the supplemental characters shows them to be largely associable with
individual QWERTY marks:

\begin{itemize}
\item
  The grave-accented vowels consonants (\textbf{á/Á, é/É, í/Í, ó/Ó, ú/Ú,
  ç/Ç, ñ/Ñ}) are accented versions of QWERTY marks.
\item
  The sentence-initial punctuation marks (\textbf{¡¿}) are inverted
  versions of QWERTY marks (\textbf{!?}).
\item
  The \emph{\textbf{guillemet}} (\textbf{« »}) are visually similar to
  the greater/less than marks (\textbf{\textless{} \textgreater{}}).
\item
  The \emph{\textbf{pesata}} sign (\textbf{₧}) starts with a capital P.
\item
  The \emph{\textbf{euro}} sign (\textbf{€}) is a currency sign
  semantically associable to the dollar sign (\textbf{\$}).
\end{itemize}

The remaining character - diaresis-accented u/U (\textbf{ü/Ü}) - is not
used in Spanish lexemes but occurs frequently enough in borrow-words to
merit inclusion in µ\textbf{Letras}. µ\textbf{Letras} associates
\textbf{v/V} keyboard marks with \textbf{ü/Ü} to distinguish it
`graphemically' from grave-accented \textbf{ú/Ú} which is associated
with the unmarked vowel keys \textbf{u/U} like other Spanish vowels
(Figure 20).

\begin{longtable}[]{@{}
  >{\raggedright\arraybackslash}p{(\columnwidth - 10\tabcolsep) * \real{0.1608}}
  >{\raggedright\arraybackslash}p{(\columnwidth - 10\tabcolsep) * \real{0.1566}}
  >{\raggedright\arraybackslash}p{(\columnwidth - 10\tabcolsep) * \real{0.1608}}
  >{\raggedright\arraybackslash}p{(\columnwidth - 10\tabcolsep) * \real{0.1659}}
  >{\raggedright\arraybackslash}p{(\columnwidth - 10\tabcolsep) * \real{0.1608}}
  >{\raggedright\arraybackslash}p{(\columnwidth - 10\tabcolsep) * \real{0.1952}}@{}}
\toprule
\begin{minipage}[b]{\linewidth}\raggedright
\end{minipage} & \begin{minipage}[b]{\linewidth}\raggedright
\end{minipage} & \begin{minipage}[b]{\linewidth}\raggedright
\textbf{ñ}
\end{minipage} & \begin{minipage}[b]{\linewidth}\raggedright
\textbf{Ç}
\end{minipage} & \begin{minipage}[b]{\linewidth}\raggedright
\end{minipage} & \begin{minipage}[b]{\linewidth}\raggedright
\end{minipage} \\
\midrule
\endhead
& & \textbf{Ñ} & \textbf{Ç} & & \\
\textbf{á} & \textbf{é} & \textbf{í} & \textbf{ó} & \textbf{ú} &
\textbf{Ü} \\
\textbf{Á} & \textbf{É} & \textbf{Í} & \textbf{Ó} & \textbf{Ú} &
\textbf{Ü} \\
\textbf{¡} & \textbf{¿} & \textbf{«} & \textbf{»} & \textbf{€} & ₧ \\
\bottomrule
\end{longtable}

Figure Non-QWERTY Spanish
Letters, Symbols, and Punctuation

\begin{longtable}[]{@{}
  >{\raggedright\arraybackslash}p{(\columnwidth - 10\tabcolsep) * \real{0.1665}}
  >{\raggedright\arraybackslash}p{(\columnwidth - 10\tabcolsep) * \real{0.1622}}
  >{\raggedright\arraybackslash}p{(\columnwidth - 10\tabcolsep) * \real{0.1665}}
  >{\raggedright\arraybackslash}p{(\columnwidth - 10\tabcolsep) * \real{0.1718}}
  >{\raggedright\arraybackslash}p{(\columnwidth - 10\tabcolsep) * \real{0.1665}}
  >{\raggedright\arraybackslash}p{(\columnwidth - 10\tabcolsep) * \real{0.1665}}@{}}
\toprule
\begin{minipage}[b]{\linewidth}\raggedright
\end{minipage} & \begin{minipage}[b]{\linewidth}\raggedright
\end{minipage} & \begin{minipage}[b]{\linewidth}\raggedright
\textbf{n}
\end{minipage} & \begin{minipage}[b]{\linewidth}\raggedright
\textbf{c}
\end{minipage} & \begin{minipage}[b]{\linewidth}\raggedright
\end{minipage} & \begin{minipage}[b]{\linewidth}\raggedright
\end{minipage} \\
\midrule
\endhead
& & \textbf{N} & \textbf{C} & & \\
\textbf{a} & \textbf{E} & \textbf{I} & \textbf{O} & \textbf{U} &
\textbf{V} \\
\textbf{A} & \textbf{E} & \textbf{I} & \textbf{O} & \textbf{U} &
\textbf{V} \\
\textbf{!} & \textbf{?} & \textbf{\textless{}} & \textbf{\textgreater{}}
& \textbf{\$} & \textbf{P} \\
\bottomrule
\end{longtable}

Figure QWERTY Marks Corresponding
to Figure 19 Characters

The strategy for solving the problem of accessing
\textbf{\emph{µ}Letras} is to make them accessible via alternative
actuations (see section 2.5.1).

Device specific

Display optional

Variety of Alternate Actuations

\section{Alternative Gestures for Common
Devices}\label{alternative-gestures-for-common-devices}

\section{Keyboard Alternative
Gestures}\label{keyboard-alternative-gestures}

\section{Mouse Alternative
Gestures}\label{mouse-alternative-gestures}

\section{Alternative Gestures for Uncommon
Devices}\label{alternative-gestures-for-uncommon-devices}

\section{Visual Preview and Feedback Associated to
Gestures}\label{visual-preview-and-feedback-associated-to-gestures}

\section{Simultaneously Presenting Character and Phrase Input
Choice}\label{simultaneously-presenting-character-and-phrase-input-choice}

e

\begin{enumerate}
\def\labelenumi{\arabic{enumi}.}
\setcounter{enumi}{4}
\tightlist
\item
\end{enumerate}
